\begin{helper}[The integral representation of a coordinate increment]
Let $n\in\N$, let $\Omega\subseteq\R^n$, let $f\colon\R^n\to\R$, let
$g\colon\R^n\to\R^n$, let $K\subseteq\R^n$, let $k$ be an index in
$\{1,\dots,n\}$ and let $d_0\in\R$.  Throughout, $e_k$ is the $k$-th standard
basis vector, so $x+se_k$ translates $x$ by $s$ in the $k$-th coordinate
direction; $K_r$ denotes the closed $r$-collar
$\{x:\operatorname{dist}(x,K)\le r\}$ of $K$; $\operatorname{supp}$ is the
\emph{closed} support; $I_h$ is the interval between $0$ and $h$; and $g^k$ is
the $k$-th component of $g$.  Assume
\begin{enumerate}
\item[(H1)] $\Omega$ is open and measurable, $K$ is measurable, and $0<d_0$
  with $K_{3d_0}\subseteq\Omega$;
\item[(H2)] the collar is stable under small coordinate shifts: for every $x$
  in $K_{d_0}$ and every $s$ with $|s|\le d_0$, the point $x+se_k$ lies in
  $K_{2d_0}$;
\item[(H3)] $g$ is a weak gradient of $f$ on $\Omega$
  (\noderef{HasWeakGradientOn}).
\end{enumerate}
Let $F_h,G_h\colon\R^n\to\R$ agree with $f$ and with $g^k$ respectively on the
collar $K_{3d_0}$, let both be integrable on $\R^n$, and let
$G_h'\colon\R^n\to\R$ be measurable and integrable with
\[
(x,s)\mapsto G_h(x+se_k)\ =\ (x,s)\mapsto G_h'(x+se_k)
\qquad\text{almost everywhere on }\R^n\times\R .
\]
Then for every $h$ with $0<|h|<d_0$ and for almost every $x\in K$, the function
$s\mapsto G_h(x+se_k)$ is interval-integrable on $I_h$ and
\[
F_h(x+he_k)-F_h(x)=\int_0^hG_h(x+se_k)\,ds .
\]
The statement is given for $F_h,G_h$ agreeing with $f,g^k$ on the collar, rather
than for $f$ and $g^k$ themselves, because \noderef{HasWeakGradientOn} controls
$g$ only through an identity against test functions compactly supported in
$\Omega$ and says nothing about $f$ off $\Omega$: the consumer applies this with
the zero extensions $F_h=\mathbf1_{K_{3d_0}}f$ and
$G_h=\mathbf1_{K_{3d_0}}g^k$, which are integrable on all of $\R^n$ while $f$
need not be.  The auxiliary $G_h'$ is a measurable representative, carried as a
hypothesis because the Fubini step below needs a genuinely measurable integrand
on the product.  This is a $W^{1,q}$ statement with no exponent in it; nothing
here is an $L^2$ claim.
\end{helper}
\begin{proof}
Fix $h$ with $0<|h|<d_0$, and write $K_1=K_{d_0}$ and $K_2=K_{2d_0}$, so
$K\subseteq K_1$ and, by (H2), $x+se_k\in K_2$ whenever $x\in K_1$ and
$|s|\le|h|<d_0$.  All three collars lie inside $\Omega$ by (H1).

\emph{Step 1: the identity for mollifications.}  Let $\rho_\delta$ be a standard
mollifier at scale $\delta$ (Preamble) and put $f_\delta=F_h*\rho_\delta$,
$(g^k)_\delta=G_h*\rho_\delta$; both are defined because $F_h$ and $G_h$ are
integrable.  For $0<\delta<d_0$ and $x\in K_1$, the function
$\Phi(s)=f_\delta(x+se_k)$ is differentiable in $s$ with
$\Phi'(s)=D_kf_\delta(x+se_k)$, since $f_\delta$ is $C^\infty$ (Preamble:
convolution with a smooth compactly supported kernel) and $s\mapsto x+se_k$ is
affine.  For $|s|\le|h|$ the point $x+se_k$ lies in $K_2$, whose $\delta$-collar
still lies in $\Omega$, so $y\mapsto\rho_\delta(x+se_k-y)$ is an admissible test
function for (H3) and the weak-gradient identity gives
$D_kf_\delta=(g^k)_\delta$ there --- here is the only use of (H3), and the only
use of the containment $K_{3d_0}\subseteq\Omega$, which is what leaves room for
both the $2d_0$-shift and the $\delta$-mollification.  That derivative is
continuous, hence interval-integrable on $I_h$, so the fundamental theorem of
calculus (Preamble) yields
\[
f_\delta(x+he_k)-f_\delta(x)=\int_0^h(g^k)_\delta(x+se_k)\,ds
\qquad(x\in K_1,\ 0<\delta<d_0).
\tag{7}
\]

\emph{Step 2: both sides converge in $L^1(K)$.}  Write
$U_\delta(x)=f_\delta(x+he_k)-f_\delta(x)$,
$V_\delta(x)=\int_0^h(g^k)_\delta(x+se_k)\,ds$, and let $U,V$ be the two sides
of the assertion, with $F_h,G_h$ in place of $f_\delta,(g^k)_\delta$.

For $U$: by the $L^1$ triangle inequality (Minkowski, Preamble),
$\int_K|U_\delta-U|\le\int_K|f_\delta(\cdot+he_k)-F_h(\cdot+he_k)|+\int_K|f_\delta-F_h|$.
Substituting $y=x+he_k$ in the first integral and using the translation
invariance of Lebesgue measure (Preamble) turns it into
$\int_{K+he_k}|f_\delta-F_h|\le\int_{K_2}|f_\delta-F_h|$, by (H2).  Both terms
tend to $0$ because mollifications converge in $L^1$ on compact sets
(Preamble).

For $V$: the function $(x,s)\mapsto|(g^k)_\delta-G_h|(x+se_k)$ is nonnegative
and measurable on $\R^n\times\R$.  Measurability is where the hypothesis on
$G_h'$ is used: $(y,s)\mapsto|(g^k)_\delta-G_h'|(y)$ is measurable, being
independent of $s$; the shear $T(x,s)=(x+se_k,s)$ is a linear automorphism of
$\R^{n+1}$ of determinant $1$, hence preserves Lebesgue measure and
measurability (Preamble); and $G_h=G_h'$ almost everywhere on the product along
the shift, so the composition is measurable.  Therefore, using
$\bigl|\int_0^hF\bigr|\le\int_{I_h}|F|$ and then Tonelli (Preamble),
\[
\int_K|V_\delta-V|
\ \le\ \int_K\!\!\int_{I_h}\bigl|(g^k)_\delta-G_h\bigr|(x+se_k)\,ds\,dx
= \int_{I_h}\!\!\int_K\bigl|(g^k)_\delta-G_h\bigr|(x+se_k)\,dx\,ds .
\]
For each fixed $s\in I_h$ the inner integral is
$\int_{K+se_k}|(g^k)_\delta-G_h|\le\int_{K_2}|(g^k)_\delta-G_h|$ by translation
invariance and (H2), so the whole expression is at most
$|h|\int_{K_2}|(g^k)_\delta-G_h|\to0$.  The same computation with
$(g^k)_\delta$ replaced by $0$ gives $\int_K|V|\le|h|\int_{K_2}|G_h|<\infty$, so
$V$ is finite almost everywhere on $K$ and lies in $L^1(K)$; likewise
$U\in L^1(K)$ because $F_h$ is integrable.  In particular, for almost every
$x\in K$ the inner integral $\int_{I_h}|G_h(x+se_k)|\,ds$ is finite, which is
the asserted interval integrability of $s\mapsto G_h(x+se_k)$.

\emph{Step 3: conclusion.}  By (7), $U_\delta=V_\delta$ everywhere on $K$, so
\[
\int_K|U-V|\ \le\ \int_K|U-U_\delta|+\int_K|V_\delta-V|\ \longrightarrow\ 0
\qquad(\delta\downarrow0),
\]
while the left-hand side does not depend on $\delta$; hence it is $0$.  Since
$|U-V|\ge0$ and is integrable on $K$, a nonnegative integrable function with
vanishing integral vanishes almost everywhere (Preamble), so $U=V$ almost
everywhere on $K$, which is the assertion.
\end{proof}
